\chapter{Sets}
\label{ch:sets}

We now fix the world that logic talks about. It contains a single kind of object,
\emph{sets}, and a single basic relation between them, \emph{membership}.

\section{Membership}

\begin{definition}[Membership]
\label{def:membership}
For objects $x$ and $A$, $x \in A$ is a proposition: \enquote{$x$ is an element of
$A$}. We write $x \notin A$ for $\lnot (x \in A)$.
\register{sets}{$x \in A$}{$x$ is an element of $A$}
\register{sets}{$x \notin A$}{$\lnot (x \in A)$}
\end{definition}

There is nothing else: no numbers, no points, no symbols. Every mathematical
object in this book, including numbers, pairs, maps and later functions and
vectors, will be \emph{a set built in a specific way}.

What we may assume about sets is fixed by the axioms of \emph{Zermelo--Fraenkel
set theory} (ZF). They are introduced one at a time, each where it is first
needed.

\section{Equality of sets}

\begin{axiom}[Extensionality]
\label{ax:extensionality}
Two sets with the same elements are equal:
\[
  \forall A\, \forall B\; \bigl( \forall x\, (x \in A \Leftrightarrow x \in B)
  \Rightarrow A = B \bigr).
\]
\end{axiom}

So a set is \emph{nothing more} than its elements. Order and repetition don't
count: $\set{a, b} = \set{b, a} = \set{a, a, b}$, once these are defined
(\cref{ax:pairing}).

\begin{definition}[Subset]
\label{def:subset}
$A \subseteq B \defiff \forall x\, (x \in A \Rightarrow x \in B)$. We say $A$ is a
\emph{subset} of $B$. A \emph{proper subset}, $A \subsetneq B$, is one with $A
\subseteq B$ and $A \neq B$.
\register{sets}{$A \subseteq B$}{$A$ is a subset of $B$}
\register{sets}{$A \subsetneq B$}{$A \subseteq B$ and $A \neq B$}
\end{definition}

\begin{theorem}
\label{thm:subset}
For all sets $A, B, C$:
\begin{enumerate}
  \item $A \subseteq A$;
  \item if $A \subseteq B$ and $B \subseteq C$, then $A \subseteq C$;
  \item $A = B$ if and only if $A \subseteq B$ and $B \subseteq A$.
\end{enumerate}
\end{theorem}

\begin{proof}
(1) For any $x$, $x \in A \Rightarrow x \in A$ is a tautology.

(2) Let $x \in A$. Then $x \in B$ because $A \subseteq B$, and then $x \in C$
because $B \subseteq C$ (chaining, \cref{thm:inference}).

(3) If $A = B$, then (1) gives both inclusions by substitution. Conversely, the
two inclusions together say $\forall x\, (x \in A \Leftrightarrow x \in B)$
(\cref{thm:logic-laws}), so $A = B$ by \cref{ax:extensionality}.
\end{proof}

Part (3) is \emph{the} standard way to prove two sets equal: show each is
contained in the other.

\section{The empty set}

\begin{axiom}[Empty set]
\label{ax:empty}
There is a set with no elements: $\exists E\; \forall x\; x \notin E$.
\end{axiom}

\begin{theorem}
\label{thm:empty-unique}
There is exactly one set with no elements.
\end{theorem}

\begin{proof}
Existence is \cref{ax:empty}. If $E$ and $E'$ both have no elements, then $x \in
E \Leftrightarrow x \in E'$ holds for every $x$, since both sides are false. So
$E = E'$ by \cref{ax:extensionality}.
\end{proof}

\begin{definition}[Empty set]
\label{def:emptyset}
This unique set is the \emph{empty set}, written $\varnothing$.
\register{sets}{$\varnothing$}{the empty set}
\end{definition}

\begin{corollary}
\label{cor:empty-subset}
$\varnothing \subseteq A$ for every set $A$.
\end{corollary}

\begin{proof}
$x \in \varnothing \Rightarrow x \in A$ is true for every $x$, because its
hypothesis is always false (\cref{rem:vacuous}).
\end{proof}

\section{Forming subsets}

\begin{axiom}[Separation]
\label{ax:separation}
For every set $A$ and every predicate $P(x)$ there is a set whose elements are
exactly the elements of $A$ that satisfy $P$:
\[
  \exists B\; \forall x\; \bigl( x \in B \Leftrightarrow (x \in A \land P(x))
  \bigr).
\]
By \cref{ax:extensionality}, $B$ is unique. We write it $\setof{x \in A}{P(x)}$.
\register{sets}{$\setof{x \in A}{P(x)}$}{the set of elements of $A$ satisfying
$P$}
\end{axiom}

This is really one axiom for each predicate $P$: an \emph{axiom schema}. The
predicate may mention other sets as parameters.

Separation only carves subsets out of a set we \emph{already have}. That
restriction is essential:

\begin{theorem}[Russell]
\label{thm:russell}
There is no set that contains every set.
\end{theorem}

\begin{proof}
Suppose $V$ were such a set. By \cref{ax:separation}, form $R = \setof{x \in
V}{x \notin x}$. Since $R$ is a set, $R \in V$. Then
\[
  R \in R \;\Leftrightarrow\; (R \in V \land R \notin R) \;\Leftrightarrow\; R
  \notin R,
\]
a proposition equivalent to its own negation. That is a contradiction.
\end{proof}

So \enquote{the collection of all sets} is not a set. Collections like that are
sometimes called \emph{classes}. We only use them as a way of speaking, never as
objects.

\section{Building larger sets}

\begin{axiom}[Pairing]
\label{ax:pairing}
For all $a, b$ there is a set whose elements are exactly $a$ and $b$. By
\cref{ax:extensionality} it is unique, written $\set{a, b}$. We write $\set{a}
\coloneqq \set{a, a}$, the \emph{singleton} of $a$.
\register{sets}{$\set{a, b}$}{the set with exactly the elements $a$ and $b$}
\register{sets}{$\set{a}$}{the set with exactly the element $a$}
\end{axiom}

So $x \in \set{a, b} \Leftrightarrow (x = a \lor x = b)$, and $x \in \set{a}
\Leftrightarrow x = a$.

\begin{axiom}[Union]
\label{ax:union}
For every set $\mathcal{F}$ there is a set whose elements are exactly the elements
of the elements of $\mathcal{F}$. It is unique, written
\[
  \bigcup \mathcal{F}, \qquad x \in \bigcup \mathcal{F} \;\Leftrightarrow\;
  \exists F \in \mathcal{F}\; x \in F.
\]
\register{sets}{$\bigcup \mathcal{F}$}{the set of elements of elements of
$\mathcal{F}$}
\end{axiom}

\begin{definition}[Union, intersection, difference]
\label{def:set-operations}
For sets $A, B$:
\[
  A \cup B \coloneqq \bigcup \set{A, B}, \qquad
  A \cap B \coloneqq \setof{x \in A}{x \in B}, \qquad
  A \setminus B \coloneqq \setof{x \in A}{x \notin B}.
\]
For a \emph{nonempty} set $\mathcal{F}$, pick any $F_0 \in \mathcal{F}$ and let
\[
  \bigcap \mathcal{F} \coloneqq \setof{x \in F_0}{\forall F \in \mathcal{F}\; x
  \in F}.
\]
This doesn't depend on the choice of $F_0$: any element of every $F \in
\mathcal{F}$ is in particular in $F_0$. Sets $A, B$ are \emph{disjoint} if $A
\cap B = \varnothing$. For three elements, $\set{a, b, c} \coloneqq \set{a, b}
\cup \set{c}$, and so on.
\register{sets}{$A \cup B$}{union of $A$ and $B$}
\register{sets}{$A \cap B$}{intersection of $A$ and $B$}
\register{sets}{$A \setminus B$}{elements of $A$ not in $B$}
\register{sets}{$\bigcap \mathcal{F}$}{elements of every element of
$\mathcal{F}$, for $\mathcal{F} \neq \varnothing$}
\register{sets}{$\set{a, b, c}$, \dots}{$\set{a, b} \cup \set{c}$, and so on}
\end{definition}

Unwinding the definitions:
$x \in A \cup B \Leftrightarrow x \in A \lor x \in B$;
$x \in A \cap B \Leftrightarrow x \in A \land x \in B$;
$x \in A \setminus B \Leftrightarrow x \in A \land x \notin B$.

\begin{remark}
$\bigcap \varnothing$ is left undefined on purpose. \enquote{The set of all $x$ in
every $F \in \varnothing$} would contain every set, which \cref{thm:russell}
forbids.
\end{remark}

\begin{theorem}[Algebra of sets]
\label{thm:set-algebra}
For all sets $A, B, C$:
\begin{center}
\begin{tabular}{@{}ll@{}}
  \toprule
  Law & Statement \\
  \midrule
  Commutativity & $A \cup B = B \cup A$, \quad $A \cap B = B \cap A$ \\
  Associativity & $(A \cup B) \cup C = A \cup (B \cup C)$, \quad same for $\cap$ \\
  Distributivity & $A \cap (B \cup C) = (A \cap B) \cup (A \cap C)$ \\
    & $A \cup (B \cap C) = (A \cup B) \cap (A \cup C)$ \\
  De Morgan & $C \setminus (A \cup B) = (C \setminus A) \cap (C \setminus B)$ \\
    & $C \setminus (A \cap B) = (C \setminus A) \cup (C \setminus B)$ \\
  Absorption & $A \cup (A \cap B) = A$, \quad $A \cap (A \cup B) = A$ \\
  Inclusion & $A \subseteq B \Leftrightarrow A \cup B = B \Leftrightarrow A \cap B = A$ \\
  \bottomrule
\end{tabular}
\end{center}
\end{theorem}

\begin{proof}
Each law turns into the matching law of logic (\cref{thm:logic-laws}) once we
unwind membership. For example, for the first distributive law and any $x$:
\begin{align*}
  x \in A \cap (B \cup C)
  &\Leftrightarrow x \in A \land (x \in B \lor x \in C) \\
  &\Leftrightarrow (x \in A \land x \in B) \lor (x \in A \land x \in C) \\
  &\Leftrightarrow x \in (A \cap B) \cup (A \cap C).
\end{align*}
So the two sets have the same elements and are equal by
\cref{ax:extensionality}. For the first De Morgan law: $x \in C \land \lnot(x \in
A \lor x \in B) \Leftrightarrow (x \in C \land x \notin A) \land (x \in C \land x
\notin B)$. The remaining laws work the same way.
\end{proof}

\begin{axiom}[Power set]
\label{ax:power}
For every set $A$ there is a set whose elements are exactly the subsets of $A$.
It is unique, written $\Pow A$:
\[
  X \in \Pow A \;\Leftrightarrow\; X \subseteq A.
\]
\register{sets}{$\Pow A$}{the set of all subsets of $A$}
\end{axiom}

For example, $\Pow \varnothing = \set{\varnothing}$ and $\Pow \set{\varnothing}
= \set{\varnothing, \set{\varnothing}}$. Note $\varnothing \neq
\set{\varnothing}$: the first has no elements, the second has one.

So from nothing, the empty set, we already get infinitely many different sets:
$\varnothing$, $\set{\varnothing}$, $\set{\set{\varnothing}}$, $\set{\varnothing,
\set{\varnothing}}$, \dots{} In \cref{ch:naturals} some of them become the
natural numbers.

\section{The remaining axioms}

ZF has four more axioms:
\begin{itemize}
  \item \textbf{Infinity} guarantees a set containing all natural numbers. It is
    introduced in \cref{ch:naturals}, where it is needed.
  \item \textbf{Replacement} says the image of a set under a definable rule is a
    set. Nothing in the chapters on numbers needs it. It will be introduced when
    it is.
  \item \textbf{Foundation} says every nonempty set $A$ has an element disjoint
    from $A$. It rules out sets like $x \in x$. We won't need it:
    \cref{lem:shape} proves $n \notin n$ for natural numbers directly.
  \item \textbf{Choice} (which turns ZF into ZFC) is not used in the number
    constructions. Whenever it is used later, that will be stated.
\end{itemize}

\begin{exercises}
  \item Prove $A \setminus (A \setminus B) = A \cap B$.
  \item Prove $A \subseteq B \Rightarrow \Pow A \subseteq \Pow B$.
  \item Write out the elements of $\Pow \set{\varnothing, \set{\varnothing}}$. How
    many are there?
  \item Show that for every set $A$ there is a set not in $A$. (Hint: adapt the
    proof of \cref{thm:russell}.)
\end{exercises}
